\documentclass[../../main.tex]{subfiles}
\begin{document}


{\bf [解]}
$a,b>0$の元で考える．$f(x)\in\mathbb Z$の定義から
\begin{align}
 x\le f(x)<x+1 \label{eq:1}
\end{align}
である．$x=ax-7,bx+3$として
\begin{align}
 ax-7\le f(ax-7)<ax-6,\qquad bx+3\le f(bx+3)<bx+4 \label{eq:2}
\end{align}
である．

さて，新しく
\begin{align}
 g(x) = \dfrac{1}{f(ax-7)} - \dfrac{1}{f(bx+3)} \label{eq:3}
\end{align}
とおく．\cref{eq:2}の各項は$x$が十分大きければ正だから
\begin{align}
 & \dfrac{1}{ax-6}-\dfrac{1}{bx+3}\ <\ g(x) <\ \dfrac{1}{ax-7}-\dfrac{1}{bx+4} \\
 & \dfrac{(b-a)x+9}{(ax-6)(bx+3)}\  <\ g(x) <\ \dfrac{(b-a)x+11}{(ax-7)(bx+4)}
\end{align}
である．
$x\to\infty$を考えるので$x>0$としてよく，両辺に$x^c (>0)$をかけると
\begin{align}
 \dfrac{(b-a)x+9}{(ax-6)(bx+3)}x^c\ <\ g(x)x^c\ <\ \dfrac{(b-a)x+11}{(ax-7)(bx+4)}x^c \label{eq:4}
\end{align}
である．この両辺の次数は$a=b$か否かで異なるので場合分けして考える．

\subparagraph{(A) $a\neq b$の時．} 

\cref{eq:4}の分子・分母を$x$で割ると
\begin{align}
\dfrac{b-a+\dfrac{9}{x}}{\bigl(a-\dfrac{6}{x}\bigr)\bigl(b+\dfrac{3}{x}\bigr)}x^{c-1}\ \le\ g(x)\,x^{c}\ <\ \dfrac{b-a+\dfrac{11}{x}}{\bigl(a-\dfrac{7}{x}\bigr)\bigl(b+\dfrac{4}{x}\bigr)}x^{c-1}
\end{align}
であり，両端の係数はともに$x\to\infty$で$\dfrac{b-a}{ab}$（$\neq 0$）に収束する．
従って$x^{c-1}$の収束・発散によって挟み撃ちの定理より$g(x)x^{c}$の収束値は以下の用に変化する
\begin{align}
 \lim_{x\to\infty}g(x)x^{c} =
 \begin{cases}
  0 & c<1 \\
  \dfrac{b-a}{ab} & c=1 \\
  \infty & c > 1
 \end{cases}
\end{align}
従って$g(x)x^{c}$が収束する$c$の最大値は$c=1$で，その時の収束値は$\dfrac{b-a}{ab}$である．


\subparagraph{(B) $a=b$の時．} 

この時は$(bx+3)=(ax-7)+10$である．
天井関数の性質より$f(y+10)=f(y)+10$だから，$f(ax+3)=f(ax-7)+10$であり，
\begin{align}
 g(x)=\dfrac{1}{f(ax-7)}-\dfrac{1}{f(ax-7)+10}=\dfrac{10}{f(ax-7)(f(ax-7)+10)}
\end{align}
と書ける．
\cref{eq:2}より$ax-7 \le f(ax-7) < ax-6$である．$x$が十分大きければこれらは全て正だから
\begin{align}
 \dfrac{10}{(ax-6)(ax+4)} < g(x) \le \dfrac{10}{(ax-7)(ax+3)}
\end{align}
である．$x^{c} (>0)$をかけて
\begin{align}
 \dfrac{10x^{c}}{(ax-6)(ax+4)} < g(x)x^c \le \dfrac{10x^{c}}{(ax-7)(ax+3)}
\end{align}
である．両辺の分母分子を$x^2$で割ると
\begin{align}
 \dfrac{10x^{c-2}}{(a-6/x)(a+4/x)} < g(x)x^c \le \dfrac{10x^{c-2}}{(a-7/x)(a+3/x)}
\end{align}
である．両辺の係数は$x\to\infty$でいずれも$\dfrac{10}{c^2}$に収束するから，挟み撃ちの定理によって，$c$の値に応じて以下のように収束値は変化する
\begin{align}
 \lim_{x\to\infty}g(x)x^{c} =
 \begin{cases}
  0 & c<2 \\
  \dfrac{10}{a^2} & c=2 \\
  \infty & c > 2
 \end{cases}
\end{align}
従って$g(x)x^{c}$が収束する$c$の最大値は$c=2$で，その時の収束値は$\dfrac{10}{a^2}$である．
\casesend

以上の議論より，求める答えは
\begin{itemize}
 \item[$a\neq b$の時] $c$の最大値は$1$，極限値は$\dfrac{b-a}{ab}$．
 \item[$a=b$の時] $c$の最大値は$2$，極限値は$\dfrac{10}{a^2}$．
\end{itemize}
である．


{\bf [解説]}
$a=b$の時は不等式の評価をより厳しくできるというのが肝である．


\newpage

\end{document}
